\scanpage{91}
de sorte que (4) s'écrit aussi
\[
\deleted{\operatorname{Tr}([u,v]w)\omega}\quad 2u\wedge v\wedge w=\operatorname{Tr}([u,v]w)\omega_{\unsure{\mathcal M|M}}\inserted{{}=\operatorname{Tr}(uvw-vuw)\omega_{\mathcal M}}\tag*{(7)}
\]
\deleted{ou encore} ou encore, compte tenu que $\operatorname{Tr}(vuw)=-\operatorname{Tr}uvw$
\[
u\wedge v\wedge w=\operatorname{Tr}(uvw)\omega_{\mathcal M}\tag*{(8)}
\]
% slanted
\margin{\uline{NB} Ceci détermine $\omega_{\mathcal M}$ \unsure{indépendamm\textsuperscript{t}|indépendamment|indépendant} de la cond. $2$ inv. sur $S$\ldots}
pour $u,v,w\in\Gamma(-,\unsure{\gamma|\mathscr V}\mathcal M)$. La loi $[\, ,\, ]$ sur $E$ est transportée de la loi $[\, ,\, ]$ sur $\gamma\mathcal M$, i.e. est définie par la loi
\[
2x\wedge y\wedge z=\varphi([x,y],z)\omega\qquad(x,y,z\in\Gamma(-,E))\tag*{(8)}
\]
% slanted
\margin{cf remarque marginale page \unsure{précédent|précédente}}
où $\varphi$ est la forme bil. sym. associée à $q$, i.e. définie par
\[
q(x+y)=q(x)+q(y)+\varphi(x,y).\tag*{(9)}
\]

Nous nous intéressons à la \uline{quadrique affine}
\[
\deleted{\unsure{Q}}\Sigma_q\inserted{{}\subset\check{\mathbb V}(E)}\quad\text{d'équation}\quad q(x)=-1\tag*{(10)}
\]
qui \deleted{correspond à la quadrique} \inserted{est can. isom. à la quadrique} affine
\[
\begin{aligned}\Sigma_{\mathcal M}\subset\check{\mathbb V}(\underbrace{\gamma\mathcal M}_{\unsure{\mathfrak g|V|\mathcal V}})&\quad\text{d'équation}\quad q_{\mathcal M}(u)=-1\quad\text{i.e.}\\&\det u=\deleted{\unsure{+}}1\end{aligned}\tag*{(11)}
\]
\margin{$\Sigma_{\mathcal M}$ est formé des $u$ \unsure{st\textsuperscript{t}|satisf\textsuperscript{t}|tels que} \deleted{qui $u^2+1=0$} à l'équation de H.C. $\boxed{u^2+1=0}$}
d'Anneaux de coordonnées affines
\[
\operatorname{Sym}^{*}(\check E)/(q+1)\operatorname{Sym}^{*}(\check E)\simeq\operatorname{Sym}^{*}(\underbrace{(\gamma\mathcal M)^\vee}_{\unsure{\mathfrak g|\mathfrak g^\vee|\mathcal V}\simeq\mathcal M/\underline O_S})/(q_{\mathcal M}+1)\tag*{(12)}
\]

Nous voulons la mettre en relation avec la S-droite projective définie par la quadrique $q(x)=0$, soit \deleted{$P$} $X$, \deleted{Pour} dont la donnée,
% archive sheet number: 90
